% Introduction (Kurose 5-paragraph structure: motivation, specific problem,
% contributions, differences from prior work, roadmap).

\section{Introduction}
\label{sec:intro}

Low Earth orbit (LEO) mega-constellations have become a mainstream component of
the global Internet: hundreds to thousands of satellites at a few hundred
kilometers of altitude now deliver low-latency broadband to regions that
terrestrial infrastructure cannot economically reach~\cite{smallsats2020,
handley:ground-relays:2019}.
Their low altitude yields round-trip times competitive with, and sometimes
better than, long-haul fiber, which has driven intense study of constellation
design and performance at scale~\cite{kassing2020exploring,lai2020starperf}.
As these systems grow to thousands of nodes carrying production traffic, the
network layer faces a defining challenge: it must keep packets flowing while the
physical topology is in perpetual, rapid motion.

Because satellites orbit continuously, a LEO network is most naturally described
as a \emph{time-varying graph} in which inter-satellite and ground--satellite
links (GSLs) appear and disappear on a schedule dictated by orbital
mechanics~\cite{seong:topological-design:1995,kassing2020exploring}.
A ground station must therefore hand its uplink from a setting satellite to a
rising one at recurring, predictable instants---a GSL \emph{handover}.
Distributed link-state routing such as OSPF~\cite{sidhu:ospf:1993}, the de facto
baseline in container-based LEO emulators~\cite{lai2023starrynet}, learns each
mutation reactively through hello timeouts and link-state flooding, and
reconverges only \emph{after} the link change has completed.
During that interval, traffic is forwarded along stale or missing next-hops and
is silently black-holed.
This raises the central question of this paper: since handover times are known
in advance from orbital prediction, can routing exploit that knowledge to keep
the data plane continuously connected across topology changes?

In this paper we present \textbf{OrbitGraph}, a graph-centric, software-defined
routing scheme for LEO constellations.
OrbitGraph models each orbital snapshot as a delay-weighted graph $G_t$, computes
all-pairs shortest paths centrally with Dijkstra's algorithm, and programs the
kernel forwarding information base (FIB) of each node by pushing only the routes
that changed since the previous snapshot.
At a handover, the controller installs the post-handover FIB \emph{before} the
old GSL is torn down, turning orbital predictability into make-before-break
forwarding at the topology level.
We implement OrbitGraph on the StarryNet emulator~\cite{lai2023starrynet}, which
runs real TCP/IP stacks and an OSPF control plane in containers, and evaluate it
against that distributed baseline on Walker constellations from 27 to 102 nodes
using a continuous outage probe aligned to routing events.
Our contributions are:
\begin{itemize}
  \item a graph-centric SDN control model that unifies initialization, delay
        refresh, link failure, and GSL handover as shortest-path computations
        over a sequence of snapshots $G_t$ (Section~\ref{sec:orbitgraph});
  \item an incremental, make-before-break FIB install that programs only changed
        next-hops, reducing handover updates by roughly two orders of magnitude
        relative to a full-table push;
  \item a proactive handover mechanism that commits the new forwarding state
        while both contacts briefly coexist; and
  \item an empirical study showing that OrbitGraph yields \mbox{0\,ms} mean
        data-plane outage at GSL handover where OSPF incurs multi-second black
        holes, holds 0\% packet loss through the handover phases on every grid
        size, and keeps shortest-path recomputation in the millisecond range,
        while also delimiting---honestly---the regime (unplanned bulk failure)
        where centralized install does not yet outperform distributed
        reconvergence (Section~\ref{sec:analysis}).
\end{itemize}

OrbitGraph differs from prior LEO routing along a single, deliberate axis:
\emph{what} forwarding state to install and \emph{when} relative to link
lifetime.
Snapshot and offline schemes precompute per-state tables and load them on
board~\cite{seong:topological-design:1995}; distributed protocols optimize hop
count or convergence speed without a global view~\cite{gregory:satellite-routing:2022,
pam:opspf:2019}; and satellite SDN research has largely targeted controller
\emph{placement} rather than route computation~\cite{papa:sdn-cpp:2018,
guo:spda:2022}.
To our knowledge, none of these install real kernel forwarding state on emulated
satellite nodes and measure end-to-end availability through handover under
proactive update ordering.
A detailed comparison appears in Section~\ref{sec:related}.

The remainder of this paper is organized as follows.
Section~\ref{sec:sdn} reviews the software-defined and graph-centric networking
principles OrbitGraph builds on, and Section~\ref{sec:related} positions it
against prior work.
Section~\ref{sec:simulation} describes the StarryNet emulation platform and
Section~\ref{sec:orbitgraph} details the OrbitGraph design.
Section~\ref{sec:experiments} presents the methodology and scaling study, whose
results are analyzed in Section~\ref{sec:analysis}.
Section~\ref{sec:future} discusses limitations and future directions, and
Section~\ref{sec:conclusion} concludes.
